\documentclass[10pt,a4paper]{amsart}
\usepackage{amsmath,amssymb,amsthm,graphicx}
% Retained preamble for the Stacks Project chapter snapshot in this directory.
% The chapter's own preamble.tex is a copy of the shared Stacks Project preamble; the block
% below is sliced MECHANICALLY, line for line, from it. Nothing was added, removed or
% redefined; the wrapper class, the xr-hyper/hyperref wiring and the externaldocument list
% are omitted because the assembler supplies the document class.
%
% lines taken from preamble.tex: 3, 5, 9, 11, 12, 156, 157, 158, ... (51 lines total)
%
% package coverage restored from the chapter's own preamble (2026-10-13): xy, xy-2cell

\usepackage[all]{xy}
\xyoption{2cell}
\UseAllTwocells

\newenvironment{reference}{\comment}{\endcomment}
\newenvironment{slogan}{\comment}{\endcomment}
\newenvironment{history}{\comment}{\endcomment}
\theoremstyle{plain}
\newtheorem{theorem}[subsection]{Theorem}
\newtheorem{proposition}[subsection]{Proposition}
\newtheorem{lemma}[subsection]{Lemma}
\theoremstyle{definition}
\newtheorem{definition}[subsection]{Definition}
\newtheorem{example}[subsection]{Example}
\newtheorem{exercise}[subsection]{Exercise}
\newtheorem{situation}[subsection]{Situation}
\theoremstyle{remark}
\newtheorem{remark}[subsection]{Remark}
\newtheorem{remarks}[subsection]{Remarks}
\numberwithin{equation}{subsection}
\def\lim{\mathop{\mathrm{lim}}\nolimits}
\def\colim{\mathop{\mathrm{colim}}\nolimits}
\def\Spec{\mathop{\mathrm{Spec}}}
\def\Hom{\mathop{\mathrm{Hom}}\nolimits}
\def\Ext{\mathop{\mathrm{Ext}}\nolimits}
\def\SheafHom{\mathop{\mathcal{H}\!\mathit{om}}\nolimits}
\def\SheafExt{\mathop{\mathcal{E}\!\mathit{xt}}\nolimits}
\def\Sch{\mathit{Sch}}
\def\Mor{\mathop{\mathrm{Mor}}\nolimits}
\def\Ob{\mathop{\mathrm{Ob}}\nolimits}
\def\Sh{\mathop{\mathit{Sh}}\nolimits}
\def\NL{\mathop{N\!L}\nolimits}
\def\CH{\mathop{\mathrm{CH}}\nolimits}
\def\proetale{{pro\text{-}\acute{e}tale}}
\def\etale{{\acute{e}tale}}
\def\QCoh{\mathit{QCoh}}
\def\Ker{\mathop{\mathrm{Ker}}}
\def\Im{\mathop{\mathrm{Im}}}
\def\Coker{\mathop{\mathrm{Coker}}}
\def\Coim{\mathop{\mathrm{Coim}}}
\DeclareMathSymbol{\boxtimes}{\mathbin}{AMSa}{"02}
\def\QCohstack{\mathcal{QC}\!\mathit{oh}}
\def\Cohstack{\mathcal{C}\!\mathit{oh}}
\def\Spacesstack{\mathcal{S}\!\mathit{paces}}
\def\Quotfunctor{\mathrm{Quot}}
\def\Hilbfunctor{\mathrm{Hilb}}
\def\Curvesstack{\mathcal{C}\!\mathit{urves}}
\def\Polarizedstack{\mathcal{P}\!\mathit{olarized}}
\def\Complexesstack{\mathcal{C}\!\mathit{omplexes}}
\def\Pic{\mathop{\mathrm{Pic}}\nolimits}
\def\Picardstack{\mathcal{P}\!\mathit{ic}}
\def\Picardfunctor{\mathrm{Pic}}
\def\Deformationcategory{\mathcal{D}\!\mathit{ef}}
\usepackage{cleveref}
\crefname{originallemma}{original lemma}{original lemmas}
\Crefname{originallemma}{Original lemma}{Original lemmas}
\makeatletter
\newlabel{equation-comparison-algebraic-spaces-modules}{{equation-comparison-algebraic-spaces-modules}{source}{}{sourcepointer-1}{}}
\newlabel{equation-comparison-algebraic-spaces-modules@cref}{{[originallemma][0][]equation-comparison-algebraic-spaces-modules}{[1][0][]source}{}{}{}}
\makeatother
\crefname{originallemma}{original lemma}{original lemmas}
\Crefname{originallemma}{Original lemma}{Original lemmas}
\makeatletter
\newlabel{equation-comparison-modules}{{equation-comparison-modules}{source}{}{sourcepointer-2}{}}
\newlabel{equation-comparison-modules@cref}{{[originallemma][0][]equation-comparison-modules}{[1][0][]source}{}{}{}}
\makeatother
\crefname{originallemma}{original lemma}{original lemmas}
\Crefname{originallemma}{Original lemma}{Original lemmas}
\makeatletter
\newlabel{lemma-stackification-quasi-coherent}{{lemma-stackification-quasi-coherent}{source}{}{sourcepointer-3}{}}
\newlabel{lemma-stackification-quasi-coherent@cref}{{[originallemma][0][]lemma-stackification-quasi-coherent}{[1][0][]source}{}{}{}}
\makeatother
\crefname{originallemma}{original lemma}{original lemmas}
\Crefname{originallemma}{Original lemma}{Original lemmas}
\makeatletter
\newlabel{proposition-quasi-coherent}{{proposition-quasi-coherent}{source}{}{sourcepointer-4}{}}
\newlabel{proposition-quasi-coherent@cref}{{[originallemma][0][]proposition-quasi-coherent}{[1][0][]source}{}{}{}}
\makeatother
\crefname{originallemma}{original lemma}{original lemmas}
\Crefname{originallemma}{Original lemma}{Original lemmas}
\makeatletter
\newlabel{spaces-groupoids-definition-quotient-stack}{{044Q}{source}{}{sourcepointer-5}{}}
\newlabel{spaces-groupoids-definition-quotient-stack@cref}{{[originallemma][0][]044Q}{[1][0][]source}{}{}{}}
\makeatother
\title{Stacks, Sheaves on Algebraic Stacks: lemma-map-from-quasi-coherent}
\author{The Stacks Project Authors}
\date{stacks-project:stacks-sheaves.tex:local-snapshot-20261012; GNU FDL 1.2 or later}
\begin{document}
\maketitle
\noindent\textit{Editorial note.} Statement span is lines 2463-2504 of sources/stacks-algebraic-stacks-sheaves/pinned-original-source.tex; proof span is lines 2506-2542, the complete author proof (35 body lines).\par\medskip
\noindent\textit{Editorial note.} Navigation stubs: 5 cited-but-not-included labels are supplied as non-load-bearing pointers with their pinned tags where the snapshot tag index has them.\par\medskip
\noindent\textit{Editorial note.} extra\_packages is empty on purpose: the retained preamble already loads xy with options and a second declaration clashes in the assembler wrapper.\par\medskip
\noindent\textit{Editorial note.} Text is retained verbatim, including the author's own typos and placeholders.\par\medskip
\noindent\textit{Editorial note.} Queue index 17 of queue\_algebraic-stacks-sheaves.json (pb=1545, stmt 2463-2504, proof 2506-2542).\par\medskip
\noindent\textit{Allowed standard prerequisites (editorial).}\begin{itemize}
\item Explicit prerequisite boundary. The cited inputs are prerequisites; the argument in the retained span is not replaced by any cited result.
\end{itemize}\medskip
\ifdefined\hypertarget\hypertarget{sourcepointer-1}{}\fi
\noindent\textit{Source reference.} Original lemma equation-comparison-algebraic-spaces-modules: sources/stacks-algebraic-stacks-sheaves/pinned-original-source.tex, source lines 1721--1725. Cited by the retained statement or proof but not part of the collected spans; supplied only as a navigation pointer. The argument does not depend on its content.\par\medskip
\ifdefined\hypertarget\hypertarget{sourcepointer-2}{}\fi
\noindent\textit{Source reference.} Original lemma equation-comparison-modules: sources/stacks-algebraic-stacks-sheaves/pinned-original-source.tex, source lines 1382--1386. Cited by the retained statement or proof but not part of the collected spans; supplied only as a navigation pointer. The argument does not depend on its content.\par\medskip
\ifdefined\hypertarget\hypertarget{sourcepointer-3}{}\fi
\noindent\textit{Source reference.} Original lemma lemma-stackification-quasi-coherent: sources/stacks-algebraic-stacks-sheaves/pinned-original-source.tex, source lines 2197--2207. Cited by the retained statement or proof but not part of the collected spans; supplied only as a navigation pointer. The argument does not depend on its content.\par\medskip
\ifdefined\hypertarget\hypertarget{sourcepointer-4}{}\fi
\noindent\textit{Source reference.} Original lemma proposition-quasi-coherent: sources/stacks-algebraic-stacks-sheaves/pinned-original-source.tex, source lines 2351--2357. Cited by the retained statement or proof but not part of the collected spans; supplied only as a navigation pointer. The argument does not depend on its content.\par\medskip
\ifdefined\hypertarget\hypertarget{sourcepointer-5}{}\fi
\noindent\textit{Source reference.} Original lemma 044Q: sources/stacks-algebraic-stacks-sheaves/pinned-original-source.tex, source lines 2542--2542. Cited by the retained statement or proof but not part of the collected spans; supplied only as a navigation pointer. The argument does not depend on its content.\par\medskip
\begin{lemma}
\label{lemma-map-from-quasi-coherent}
Let $(U, R, s, t, c)$ be a groupoid in algebraic spaces over $S$.
Assume $s, t$ are flat and locally of finite presentation.
Let $\mathcal{X} = [U/R]$ be the quotient stack. Denote
$x$ the object of $\mathcal{X}$ over $U$.
Let $\mathcal{F}$ be a quasi-coherent
$\mathcal{O}_\mathcal{X}$-module, and let $\mathcal{H}$ be any object
of $\textit{Mod}(\mathcal{O}_\mathcal{X})$.
The map
$$
\Hom_{\mathcal{O}_\mathcal{X}}(\mathcal{F}, \mathcal{H})
\longrightarrow
\Hom_{\mathcal{O}_U}(x^*\mathcal{F}|_{U_\etale},
x^*\mathcal{H}|_{U_\etale}),
\quad
\phi \longmapsto x^*\phi|_{U_\etale}
$$
is injective and its image consists of exactly those
$\varphi : x^*\mathcal{F}|_{U_\etale} \to
x^*\mathcal{H}|_{U_\etale}$ which give rise to a commutative
diagram
$$
\xymatrix{
s_{small}^*(x^*\mathcal{F}|_{U_\etale})
\ar[r] \ar[d]^{s_{small}^*\varphi} &
(x \circ s)^*\mathcal{F}|_{R_\etale} =
(x \circ t)^*\mathcal{F}|_{R_\etale} &
t_{small}^*(x^*\mathcal{F}|_{U_\etale})
\ar[l] \ar[d]_{t_{small}^*\varphi} \\
s_{small}^*(x^*\mathcal{H}|_{U_\etale})
\ar[r] &
(x \circ s)^*\mathcal{H}|_{R_\etale} =
(x \circ t)^*\mathcal{H}|_{R_\etale} &
t_{small}^*(x^*\mathcal{H}|_{U_\etale})
\ar[l]
}
$$
of modules on $R_\etale$
where the horizontal arrows are the comparison maps
(\ref{equation-comparison-algebraic-spaces-modules}).
\end{lemma}

\begin{proof}
According to
Lemma \ref{lemma-stackification-quasi-coherent}
the stackification map $[U/_{\!p}R] \to [U/R]$ (see
Groupoids in Spaces, Definition
\ref{spaces-groupoids-definition-quotient-stack})
induces an equivalence of categories of quasi-coherent sheaves
and of fppf $\mathcal{O}$-modules.
Thus it suffices to prove the lemma with $\mathcal{X} = [U/_{\!p}R]$.
By Proposition \ref{proposition-quasi-coherent}
and its proof there exists a quasi-coherent module
$(\mathcal{G}, \alpha)$ on $(U, R, s, t, c)$ such that
$\mathcal{F}$ is given by the rule
$\mathcal{F}(T, u) = \Gamma(T, u^*\mathcal{G})$.
In particular $x^*\mathcal{F}|_{U_\etale} = \mathcal{G}$
and it is clear that the map of the statement of the
lemma is injective. Moreover, given a map
$\varphi : \mathcal{G} \to x^*\mathcal{H}|_{U_\etale}$
and given any object
$y = (T, u)$ of $[U/_{\!p}R]$ we can consider the map
$$
\mathcal{F}(y) = \Gamma(T, u^*\mathcal{G})
\xrightarrow{u_{small}^*\varphi}
\Gamma(T, u_{small}^*x^*\mathcal{H}|_{U_\etale})
\rightarrow
\Gamma(T, y^*\mathcal{H}|_{T_\etale}) = \mathcal{H}(y)
$$
where the second arrow is the comparison map
(\ref{equation-comparison-modules}) for the sheaf $\mathcal{H}$.
This assignment is compatible with the restriction mappings of the
sheaves $\mathcal{F}$ and $\mathcal{G}$ for morphisms of
$[U/_{\!p}R]$ if the cocycle condition of
the lemma is satisfied. Proof omitted. Hint: the restriction maps
of $\mathcal{F}$ are made explicit in terms of $(\mathcal{G}, \alpha)$
in the proof of
Proposition \ref{proposition-quasi-coherent}.
\end{proof}

\end{document}
